\section{Limitations}

First, the evidence is simulation-only. It supports a mechanism claim inside a trust-task testbed, not a claim about human affect, human trust, or clinical diagnosis. Human behavioral fitting or experimental data would be needed before claiming psychological validity.

Second, the primary trust task uses environment-owned partner behavior. This is appropriate for isolating the focal agent's partner-local precision mechanism, but it does not establish fully reciprocal multi-agent dynamics. AIF-vs-AIF partner interactions remain future or descriptive work until promoted into specific hypotheses.

Third, partner-specific beta has not yet been compared against a global-beta ablation. Partner specificity is an architectural premise of the current model. A global-beta control would directly test whether a shared affective precision state can explain the same partner-choice, entropy, payoff, and model-fitness results.

Fourth, H0, H2, H4, and H5 use five seeds per variant. They are useful supporting evidence but should not carry the manuscript's strongest claims. H1 and H3, each with 30 seeds per variant, should remain the empirical anchors.

Fifth, H3 is a boundary condition rather than a recovery win. Abrupt betrayal reduces affective-agent payoff in the 30-seed confirmation. The paper should state this plainly. The mechanism changes deployment and reallocation but does not produce safer return behavior under that shock regime.

Sixth, H5 payoff claims are underpowered and task-dependent. The clinical-like variants separate in precision dynamics, but diagnostic language should be avoided. The correct terms are ``clinical-like parameter perturbations'' or ``computational perturbation variants.''

Seventh, the current figures are draft analysis aids. Camera-ready figures should be redrawn as multi-panel composites from source tables, with consistent scales, uncertainty intervals, and captions that state seed counts.

\section{Future Work}

The highest-value follow-up is a global-beta ablation. The model's architectural novelty is partner-specific affective precision, but the current evidence does not directly compare partner-local beta against one shared beta tracker. A global-beta variant should share one beta posterior across partners while keeping the same $\gamma=\gamma_{\mathrm{base}}/\mathbb{E}[\beta]$ mapping. The most useful tests would be H1 reliability-versus-reward, H2 open-regime lesion, and H4 partner choice, with optional H3 stress checks.

A second useful follow-up is a focused shock-shape gradient. The current H3 sensitivity results suggest that abruptness, not generic overconfidence, drives the betrayal failure. A focused experiment could hold default affect fixed while varying the abruptness of the stance switch, observation noise, or number of pre-switch confirmations. The primary readouts should be payoff, entropy, reencounters, payoff-on-return, wrong-type rate, and low-entropy wrong-rate.

A third near-term need is figure-quality composite generation. Current PNGs are diagnostic aids. A script that reads the copied source tables and emits final multi-panel PNG/PDF figures would likely improve the manuscript more than another broad sweep.

Longer-term work includes human-data validation, model comparison against reward-averaging and belief-only baselines, a fuller variational beta implementation, richer social environments with correlated partners, and reciprocal AIF agents interacting with each other.

\section{Conclusion}

Partner-specific affective precision is a plausible computational bridge between social model fitness and social policy deployment. In the current simulations, it is behaviorally visible when policy spaces are open, tracks predictive reliability more than payoff, changes action deployment despite similar belief readouts, and alters partner choice before payoff clearly moves. Under abrupt betrayal, the same mechanism can sharpen the wrong deployment state and reduce payoff. The current evidence therefore supports a narrow but useful claim: affective precision is a partner-local model-fitness signal, not a generic reward booster. Future work should test the partner-specific architecture against global beta, validate the mechanism against human behavior, and map the shock-regime boundary conditions more precisely.
